\scanpage{103}
\hrulefill

Soit $X/K$, et soit $T$ une partie de \inserted{$\unsure{X}=$}$\unsure{X_{\mathbb C}(\mathbb C)|X(\mathbb C)}$, loc.\textsuperscript{t} fermée 1-connexe. Soit $X'$ \deleted{une variété algébrique}\inserted{$\unsure{-|\longrightarrow}X$ un} $K$-morphisme, qui sur \deleted{la partie de} un voisinage de $T$ soit fini étale. Alors $X'\vert T=X'_T$ se décompose en un nb fini de copies de $T$ -- il y a une relation d'équivalence dans $X'_T$, qu'on appelle la $T$-adjacence de points de $X'$ au dessus de $T$. On se donne \deleted{$I\subset X'_T$} $x',y'\in X'_T$ \inserted{$T$-adjacents;} si $u\in\Gamma$, on ne peut affirmer que \deleted{\unsure{ils le soient}}\inserted{$u(x'),u(y')$} soient encore $T$ adjacents.
\margin{en supposant $x',y'\in X'_{\overline K}$}
Mais s'ils le sont, on peut regarder $l^T_{y',x'}:x'\to y'$ \inserted{``$T$-horizontal''} dans \deleted{$\unsure{\widehat{\Pi}_1(X'_{\overline K},x',y')}$} $\operatorname{Hom}_{\widehat{\Pi}_1(X'_{\overline K})}(x',y')$ et son transformé, \deleted{\unsure{$T$ \illegible{1} qui passe par $x',y'$}} \deleted{et son transformé} et le \uline{comparer} \deleted{$\unsure{u(l_{y',x'})}$} à $l^T_{u(y'),u(x')}$ (qui est défini par hyp. d'horizontalité)
\[
(53)\qquad \deleted{\unsure{u}}{}^u l^T_{y',x'}=g^T_{x,y}(u)\,l^T_{u(y'),u(x')}
\]
où
\[
g^{T\deleted{\unsure{u}}}_{x,y}(u)\in\widehat{\pi}_1\left(X',\underbrace{T(u x',u y')}_{\substack{\text{strate horizontale}\\\text{de }X'_T\text{ qui contient}\\u(x'),u(y')}}\right).
\]
Ceci posé, on aura
\[
(54)\qquad \widehat f\underbrace{\bigl(g^T_{y',x'}(u)\bigr)}_{\substack{\rotatebox{-90}{$\unsure{\in}$}\\\widehat{\pi}_1(X,T)}}=g^T_{y,x}(u)
\]
\mspage{102}
